\documentclass[10pt,a4paper]{amsart}
\usepackage[margin=19mm]{geometry}
\usepackage{amsmath,amssymb,mathtools}
\usepackage[scr=rsfs,cal=euler]{mathalfa}
\usepackage{xfrac}
\usepackage[unicode,hidelinks,bookmarks=false]{hyperref}
\newcommand{\ie}{i.e.\ }
\newcommand{\cf}{cf.\ }
\newcommand{\resp}{respectively\ }
\newcommand{\Subsection}{\subsection}
\providecommand{\nobreakdash}{\nobreak\hbox{-}}
\setlength{\emergencystretch}{2em}
\multlinegap0pt
\usepackage{bm}
\usepackage{bbm}
\usepackage{dsfont}
\usepackage{xspace}
\usepackage{cancel}
\usepackage{mathtools}
\usepackage{amsthm}
\makeatletter
\@ifundefined{theorem}{\newtheorem{theorem}{Theorem}}{}
\@ifundefined{lemma}{\newtheorem{lemma}[theorem]{Lemma}}{}
\makeatother
\title[A complete solution to the biased Alon-Krivelevich-Spencer-S]{A complete solution to the biased Alon-Krivelevich-Spencer-Szab\'o criterion problem for the discrepancy game}
\author{Yaping Mao, Meiqin Wei, Gang Yang}
\date{Source: arXiv:2606.13309v2 (CC BY 4.0)}
\begin{document}
\maketitle

\noindent\emph{Source and licence.} A complete solution to the biased Alon-Krivelevich-Spencer-Szab\'o criterion problem for the discrepancy game. Version arXiv:2606.13309v2 (CC BY 4.0), distributed under the Creative Commons Attribution 4.0 International licence, as recorded in the arXiv licence metadata for this exact version and shown on the abstract page https://arxiv.org/abs/2606.13309v2. Full rights notice: licenses/arxiv-2606.13309-CC-BY-4.0.txt.\par\smallskip

\noindent\textit{Editorial scope.} Counted unit: the Theorem whose source label is \texttt{thm:pseudorandom} in the article (source0.tex lines 345--371, source label \texttt{thm:pseudorandom}), delivered together with the article's own complete original proof of it (source0.tex lines 1303--1605, one \texttt{proof} environment of 303 lines). No mathematics is written, repaired or re-derived: statement and proof are the article's own text line for line. The proof is detached from the statement in the source: the article states the result at lines 345--371 and prints its proof at lines 1303--1605, and the proof environment's own header names the statement, so the association is the article's own. Required context retained uncounted: thm:main (lines 282--306); eq:main-condition (lines 285--293); lem:local-quadratic (lines 1170--1179). Every definition, display and label of those lines is retained unchanged. Omitted from this package and not redistributed: 1--281, 294--344, 372--1169, 1180--1302, 1606--1700. Nothing inside a retained excerpt is omitted or altered: every retained body is delivered exactly as the article's lines are written, so no omission receipt of any kind accompanies this package. Citations: the retained excerpts carry no citation command at all, so no bibliography entry of the article is transcribed and no reference is redistributed. Reference adaptation: none was needed and none was made; every reference inside the delivered unit resolves inside it, so no unresolved reference or citation is produced. Printed slips kept literal: no printed word, symbol, hypothesis or number of the counted statement or of its proof was altered. Layout and class: the article's own class is \texttt{article} with packages including fontenc, inputenc, lmodern, mathrsfs, amssymb, amsmath, mathtools, amsthm, graphicx, color, float, caption, enumitem, booktabs; the wrapper is the lane's pinned \texttt{amsart} editorial wrapper at 10pt on A4, which restates the theorem-like environments the retained text needs and adds the article's own macro definitions that the retained text needs (none), with extra wrapper packages (bm, bbm, dsfont, xspace, cancel, mathtools). The delivered numbering is the unit's own. Assets: no retained span contains a figure environment, a graphics-inclusion command, a PDF-page inclusion, a table or a TikZ picture; no image file is copied into this package, the package declares no asset dependency and no image file is redistributed with it. Licence: version arXiv:2606.13309v2 of this article is distributed under the Creative Commons Attribution 4.0 International licence, as recorded in the arXiv licence metadata for this exact version and shown on the abstract page https://arxiv.org/abs/2606.13309v2. A full rights notice is delivered at licenses/arxiv-2606.13309-CC-BY-4.0.txt. Characters: every retained excerpt is pure ASCII, so the delivered engine, XeLaTeX with its default Latin Modern text font, reports no missing character.
\begin{theorem}
\label{thm:main}
Let $H=(V,\mathcal E)$ be a finite hypergraph with no empty edges. Suppose that $|V|$ is divisible by $p+q$, and that the $(p:q)$-game on $V$ is played in complete rounds with Balancer moving first in each round. For every edge $e\in\mathcal E$, let $L_e^+,L_e^-\ge0$. Assume that there exist parameters $t_e,s_e>0$ such that
\begin{equation}
\label{eq:main-condition}
        \sum_{e\in\mathcal E} e^{-t_eL_e^+}
        \bigl(R^+_{p,q}(t_e)\bigr)^{|e|}
        +
        \sum_{e\in\mathcal E} e^{-s_eL_e^-}
        \bigl(R^-_{p,q}(s_e)\bigr)^{|e|}
        \le 1.
\end{equation}
Then Balancer has a strategy guaranteeing, at the end of the game,
\[
        -L_e^-\le q|B\cap e|-p|U\cap e|\le L_e^+
        \qquad\text{for all }e\in\mathcal E.
\]
Equivalently,
\[
        -\frac{L_e^-}{p+q}\le
        |B\cap e|-\frac{p}{p+q}|e|\le
        \frac{L_e^+}{p+q}
        \qquad\text{for all }e\in\mathcal E.
\]
\end{theorem}

\begin{lemma}\label{lem:local-quadratic}
For fixed positive integers $p,q$, there exist constants
$A=A(p,q)>0$ and $\tau=\tau(p,q)>0$ such that
\[
        \log R^+_{p,q}(x)\le A x^2,
        \qquad
        \log R^-_{p,q}(x)\le A x^2
        \qquad(0\le x\le \tau).
\]
\end{lemma}

\begin{theorem}\label{thm:pseudorandom}
Fix positive integers $p,q$, and put
\[
        \rho=\frac{p}{p+q}.
\]
There exist constants $C=C(p,q)$ and $\eta_0=\eta_0(p,q)>0$ such that the
following holds. Let $0<\eta\le\eta_0$, let $n\ge3$, and suppose that
\[
        \eta^3 n\ge C\log(e/\eta).
\]
In the $(p:q)$-game played on the edge set of $K_n$, Balancer has a strategy
to build a graph $G_B$ satisfying
\[
        \left|d_{G_B}(v)-\rho(n-1)\right|
        \le
        C\sqrt{n\log n}+O_{p,q}(1)
        \qquad\text{for every }v\in[n],
\]
and, for every pair of disjoint sets $S,T\subseteq[n]$ with
$|S|,|T|\ge\eta n$,
\[
        \left|
        \frac{e_{G_B}(S,T)}{|S||T|}-\rho
        \right|
        \le \eta .
\]
\end{theorem}

\begin{proof}[Proof of Theorem~\ref{thm:pseudorandom}]
Let $A$ and $\tau$ be the constants from Lemma~\ref{lem:local-quadratic}.
We shall choose the constants in the statement in the following order.  First choose
$K=K(p,q)$ so large that
\[
        \frac{(p+q)^2K^2}{4A}\ge 3.
\]
Next choose $\eta_0=\eta_0(p,q)>0$ so small that
\[
        0<\eta_0\le \frac14
        \qquad\text{and}\qquad
        \frac{(p+q)\eta_0}{4A}\le \tau.
\]
Finally choose $C=C(p,q)$ sufficiently large.  During the proof we shall increase
$C$ several times, always depending only on $p,q$.

Let $0<\eta\le\eta_0$, and put
        $k=\lceil \eta n\rceil$.
By increasing $C$ if necessary, the assumption
$\eta^3 n\ge C\log(e/\eta)$
implies
        $\eta n\ge1$
        and hence 
        $k\le 2\eta n$.
It also implies that $n$ is sufficiently large for all estimates below.

We first prove the assertion in the complete-round formulation.  The board of the
auxiliary discrepancy game is the edge set of $K_n$, that is, $W=E(K_n)$.
A set claimed by Balancer in this auxiliary game will be regarded as the edge set
of the graph $G_B$.

We define an auxiliary hypergraph
$\mathcal H=(W,\mathcal F)$
as follows.  The first family of hyperedges consists of the stars
$F_v=\{vx:x\in[n]\setminus\{v\}\}
  (v\in[n])$.
Thus
        $|F_v|=n-1$.
The second family of hyperedges consists of bipartite edge sets
$F_{S,T}=\{xy:x\in S,\ y\in T\}$,
where $S,T\subseteq[n]$ are disjoint $k$-sets. Thus
        $|F_{S,T}|=k^2$.
We shall apply Theorem~\ref{thm:main} to this auxiliary hypergraph.

For every star $F_v$, set
\begin{equation}\label{3-5}
L_v^+=L_v^-=(p+q)K\sqrt{(n-1)\log(2n)}.
\end{equation}
We choose the same exponential parameter for the two tails:
\begin{equation}\label{3-6}
t_v=s_v=
        \frac{L_v^+}{2A(n-1)}
        =
        \frac{(p+q)K}{2A}
        \sqrt{\frac{\log(2n)}{n-1}}.
\end{equation}
After increasing $C$, the hypothesis on $\eta^3 n$ guarantees that
$n$ is large enough so that
$t_v=s_v\le\tau$
for $v\in[n]$.

By Lemma~\ref{lem:local-quadratic},
$\log R^+_{p,q}(t_v)\le A t_v^2$ and $\log R^-_{p,q}(s_v)\le A s_v^2$. 
Therefore, by \eqref{3-5} and \eqref{3-6}, each upper star summand in \eqref{eq:main-condition} is at most
\[
\begin{aligned}
        e^{-t_vL_v^+}
        \bigl(R^+_{p,q}(t_v)\bigr)^{n-1}
        &\le
        \exp\{-t_vL_v^+ + A(n-1)t_v^2\}  \\
        &=
        \exp\left\{
        -\frac{(L_v^+)^2}{4A(n-1)}
        \right\}  \\
        &=
        \exp\left\{
        -\frac{(p+q)^2K^2}{4A}\log(2n)
        \right\}.
\end{aligned}
\]
The same estimate holds for the lower star summand.  Since
\[
        \frac{(p+q)^2K^2}{4A}\ge3,
\]
the total contribution of the upper and lower tails of all stars is at most
$2n(2n)^{-3}<\frac14$
for $n\ge3$.


We now consider the bipartite hyperedges $F_{S,T}$.  For every ordered pair
of disjoint $k$-sets $S,T$, set
\begin{equation}\label{3-7}
L_{S,T}^+=L_{S,T}^-=(p+q)\frac{\eta}{2}k^2.
\end{equation}
Again choose the same exponential parameter for the two tails:
\begin{equation}\label{3-8}
t_{S,T}=s_{S,T}
        =
        \frac{L_{S,T}^+}{2A k^2}
        =
        \frac{(p+q)\eta}{4A}.
\end{equation}
The choice of $\eta_0$ gives
$t_{S,T}=s_{S,T}\le\tau$.
Thus Lemma~\ref{lem:local-quadratic} gives
$\log R^+_{p,q}(t_{S,T})\le A t_{S,T}^2$
and $\log R^-_{p,q}(s_{S,T})\le A s_{S,T}^2$.
Hence, by \eqref{3-7} and \eqref{3-8}, each upper bipartite summand is at most
\[
\begin{aligned}
        e^{-t_{S,T}L_{S,T}^+}
        \bigl(R^+_{p,q}(t_{S,T})\bigr)^{k^2}
        &\le
        \exp\{-t_{S,T}L_{S,T}^+ + A k^2t_{S,T}^2\}  \\
        &=
        \exp\left\{
        -\frac{(L_{S,T}^+)^2}{4Ak^2}
        \right\}  \\
        &=
        \exp\left\{
        -\frac{(p+q)^2}{16A}\eta^2k^2
        \right\}.
\end{aligned}
\]
The lower bipartite summand satisfies the same bound.  Put
\[
        \gamma=\frac{(p+q)^2}{16A}.
\]
The number of ordered pairs of disjoint $k$-subsets is at most
$\binom nk^2$.
Using
$\binom nk\le \left(\frac{en}{k}\right)^k$
and $k\ge \eta n$,
we get
\[
        \binom nk^2
        \le
        \left(\frac{en}{k}\right)^{2k}
        \le
        \left(\frac e\eta\right)^{2k}.
\]
Therefore the total contribution of all upper and lower bipartite tails is at most
$2\exp\left\{
        2k\log(e/\eta)-\gamma\eta^2k^2
        \right\}$.
Since $k\ge\eta n$, we have
$\eta^2k^2\ge \eta^3nk$.
Thus the exponent is at most
$2k\log(e/\eta)-\gamma\eta^3nk$.
By increasing $C$ so that
$\gamma C\ge5$,
the hypothesis
$ \eta^3 n\ge C\log(e/\eta)$
implies
$2k\log(e/\eta)-\gamma\eta^3nk
        \le
        -3k\log(e/\eta)$.
As $k\ge1$ and $\log(e/\eta)\ge1$, the total contribution of all bipartite
hyperedges is less than
$2e^{-3}<\frac14$.

Combining the star and bipartite estimates, the sum in \eqref{eq:main-condition}
for the auxiliary hypergraph is strictly smaller than $1$.  Theorem~\ref{thm:main}
therefore applies.  Hence Balancer has a strategy such that, for every vertex
$v\in[n]$,
\[
        \left|
        |B\cap F_v|-\rho |F_v|
        \right|
        \le
        \frac{L_v^+}{p+q}
        =
        K\sqrt{(n-1)\log(2n)}.
\]
Since $B$ is the edge set of $G_B$, we have
$|B\cap F_v|=d_{G_B}(v)$
and
$|F_v|=n-1$.
Thus
\[
        \left|
        d_{G_B}(v)-\rho(n-1)
        \right|
        \le
        K\sqrt{(n-1)\log(2n)}
        \qquad(v\in[n]).
\]
After increasing the constant $C=C(p,q)$ in the statement, this gives
\[
        \left|
        d_{G_B}(v)-\rho(n-1)
        \right|
        \le
        C\sqrt{n\log n}
        \qquad(v\in[n])
\]
in the complete-round setting.

Similarly, for every ordered pair of disjoint $k$-sets $S,T$, Theorem~\ref{thm:main}
gives
\[
        \left|
        |B\cap F_{S,T}|-\rho |F_{S,T}|
        \right|
        \le
        \frac{L_{S,T}^+}{p+q}
        =
        \frac{\eta}{2}k^2.
\]
Since
$|B\cap F_{S,T}|=e_{G_B}(S,T)$
and $|F_{S,T}|=k^2$,
we obtain
\[
        \left|
        e_{G_B}(S,T)-\rho k^2
        \right|
        \le
        \frac{\eta}{2}k^2
\]
for all disjoint $k$-sets $S,T$.

It remains to discuss the possible last incomplete round.  If the board size is not
divisible by $p+q$, or if the rules force a final incomplete round, then at most
$p+q-1$ final claims are outside the complete-round analysis.  Such claims can change
the degree of any vertex by at most $p+q-1$, and can change the value of
$e_{G_B}(S,T)$ for any fixed pair $S,T$ by at most $p+q-1$.  Hence the degree estimate
gets an additional $O_{p,q}(1)$ term.  For the bipartite estimates, by increasing
$C$ once more we may assume
$p+q-1\le \frac{\eta}{2}k^2$.
Indeed, since $k\ge\eta n$, we have
$ \eta k^2\ge \eta^3 n^2\ge \eta^3 n\ge C\log(e/\eta)$,
and the right hand side can be made larger than any fixed constant depending only on
$p,q$.  Therefore, in the original game, for all disjoint $k$-sets $S,T$,
$\left|
        e_{G_B}(S,T)-\rho k^2
        \right|
        \le
        \eta k^2$.

Finally we pass from $k$-sets to arbitrary larger sets.  Let $S,T\subseteq[n]$ be
disjoint sets with
        $|S|,|T|\ge\eta n$.
Since $k=\lceil\eta n\rceil$, we have $k\le |S|$ and $k\le |T|$.  Choose uniformly
random $k$-subsets
$S'\subseteq S$ and
        $T'\subseteq T$.
For every choice of $S'$ and $T'$, the sets are disjoint, and hence the estimate just
proved gives
\begin{equation}\label{3-9}
\left|
        e_{G_B}(S',T')-\rho k^2
        \right|
        \le
        \eta k^2.
\end{equation}
Now each edge between $S$ and $T$ is
chosen in $e_{G_B}(S',T')$ with probability
\[
        \frac{k}{|S|}\cdot\frac{k}{|T|}
        =
        \frac{k^2}{|S||T|}.
\]
Therefore
\[
        \mathbb E[\,e_{G_B}(S',T')]
        =
        \frac{k^2}{|S||T|}e_{G_B}(S,T).
\]
Since \eqref{3-9}
holds for every choice of the \(k\)-subsets \(S'\subseteq S\) and
\(T'\subseteq T\), it also holds after averaging over all such choices. Hence
\[
        \left|
        \mathbb E [e_{G_B}(S',T')]-\rho k^2
        \right|
        \le
        \eta k^2 .
\]
Using
\[
        \mathbb E [e_{G_B}(S',T')]
        =
        \frac{k^2}{|S||T|}e_{G_B}(S,T),
\]
we obtain
\[
        \left|
        \frac{k^2}{|S||T|}e_{G_B}(S,T)-\rho k^2
        \right|
        \le
        \eta k^2 .
\]
Dividing by $k^2$ gives
\[
        \left|
        \frac{e_{G_B}(S,T)}{|S||T|}-\rho
        \right|
        \le
        \eta.
\]
Together with the degree estimate, this proves Theorem~\ref{thm:pseudorandom}.
\end{proof}

\end{document}
